% !TeX root = ../seminar.tex
\begin{frame}{What comes next in the thesis}
\begin{itemize}\setlength{\itemsep}{9pt}
  \item<+-> \textbf{Refine the forecast.} Run on the adopted Roman footprint at finer resolution
        and test how correlated astrometric errors affect mass precision.
  \item<+-> \textbf{Test black-hole populations.} Vary the mass distribution and the velocities
        black holes receive at birth.
  \item<+-> \textbf{Infer the population.} Use mock detections, with survey selection and mass
        errors, to test what we can learn about abundance and a possible low-mass gap.
\end{itemize}
\note{The first task is to establish which forecast results are stable as we refine the sky grid
      and vary the astrometric noise assumptions. The second asks how the predicted observations
      change with different black-hole populations. The third applies those forecasts to inference
      from a realistically selected catalogue.}
\end{frame}
